\begin{frame}{(True) Population Risk \cite{tamang2025handlingoutofdistributiondatasurvey}}
    Given an input space or domain, denoted as $\mathcal X$, an approximation function, denoted as $\hat f: \mathcal X \mapsto \mathcal Y \subseteq \mathbb R$,
    and $f$ is the true function representing the relationship $\mathcal X \mapsto \mathcal Y$.
    The models repeatedly update their parameters to \cite{gulrajani2020searchlostdomaingeneralization} diminish the population risk.
    \medskip

    \noindent
    \label{def1}
    \[
    \mathcal R_{\mathrm{true}}\left(\hat f\right) \equiv \mathcal{R}_{\mathbf x \in \mathcal X}\left(\hat f\right)
    \stackrel{\mathrm{def}}{=} \mathbb E_{\mathbf x \in \mathcal X}\!\left[\ell\big(\hat f(\mathbf x), f(\mathbf x)\big)\right]
    \]
\end{frame}

\begin{frame}{Empirical Risk \cite{tamang2025handlingoutofdistributiondatasurvey} and Population Risk}
    Nonetheless, in real-world scenarios, it is almost impossible to capture all behaviours of $\mathcal X$ to train the models.
    Alternatively, the model $\hat f$ is trained on a \textbf{sub-domain} $\mathcal X^{\text{tr}} \subset \mathcal X$, and hence, the model $\hat f$ is trying to optimize the sample risk or empirical risk.
    \label{def2}
    \[
    \mathcal R_{\mathrm{emp}}\left(\hat f\right) \equiv \mathcal{R}_{\mathbf x \in \mathcal X^{\text{tr}}}\left(\hat f\right)
    \stackrel{\mathrm{def}}{=} \mathbb E_{\mathbf x \in \mathcal X^{\text{tr}}}\!\left[\ell\big(\hat f(\mathbf x), f(\mathbf x)\big)\right]
    \]

    \noindent
    The model $\hat f$ is expected to \cite{tamang2025handlingoutofdistributiondatasurvey} make the \textbf{Empirical Risk} $\mathcal R_{\text{emp}}$ approximates the \textbf{True Risk} $\mathcal R_{\text{true}}$. Formally, this is
    \[
    \mathop{\arg\min}\limits_{\hat f} \mathcal R_{\text{emp}}\left(\hat f\right) \approx \mathop{\arg\min}\limits_{\hat f} \mathcal R_{\text{true}}\left(\hat f\right)
    \]
\end{frame}